\documentclass[11pt,a4paper]{article}
\usepackage[utf8]{inputenc}
\usepackage[T1]{fontenc}
\usepackage[french]{babel}
\usepackage{amsmath, amssymb}
\usepackage{geometry}
\geometry{margin=2.5cm}
\usepackage{xcolor}
\usepackage{listingsutf8}
\usepackage[most]{tcolorbox}
\usepackage{prftree}
\usepackage{hyperref}
\usepackage{newunicodechar}

\newunicodechar{ι}{\ensuremath{\iota}}
\newunicodechar{₁}{\ensuremath{_1}}
\newunicodechar{₂}{\ensuremath{_2}}

\newtcbtheorem{remarque}{Remarque}{
  colback=blue!5!white,
  colframe=blue!75!black,
  fonttitle=\bfseries,
  breakable,
  enhanced
}{rem}

\newtcbtheorem{exercice}{Exercice}{
  colback=green!5!white,
  colframe=green!40!black,
  fonttitle=\bfseries,
  breakable,
  enhanced
}{exo}

\definecolor{backcolour}{rgb}{0.95,0.95,0.92}
\definecolor{codegreen}{rgb}{0,0.6,0}
\lstset{
  language=ML,
  inputencoding=utf8,
  backgroundcolor=\color{backcolour},
  commentstyle=\color{codegreen},
  keywordstyle=\color{blue}\bfseries,
  basicstyle=\ttfamily\small,
  breaklines=true,
  frame=single,
  rulecolor=\color{lightgray},
  showstringspaces=false,
  extendedchars=true,
  literate={ι}{{$\iota$}}1 {₁}{{$_1$}}1 {₂}{{$_2$}}1 {é}{{\'e}}1 {è}{{\`e}}1 {à}{{\`a}}1 {ù}{{\`u}}1 {ê}{{\^e}}1 {î}{{\^i}}1 {ô}{{\^o}}1 {û}{{\^u}}1 {ë}{{\"e}}1 {ï}{{\"i}}1 {ö}{{\"o}}1 {ü}{{\"u}}1
}

\makeatletter
\newcommand{\subtitle}[1]{%
  \gdef\@subtitle{#1}%
}
\let\oldmaketitle\maketitle
\renewcommand{\maketitle}{%
  \let\oldtitle\@title
  \renewcommand{\@title}{\oldtitle\\[1em]\Large\normalfont\@subtitle}%
  \oldmaketitle
}
\makeatother

\title{LEÇON 2 : Logique Propositionnelle et Constructivisme}
\subtitle{Connecteurs ET, OU et Négation}
\date{24/03/2026}
\author{Pablo Donato}

\begin{document}

\maketitle

Après avoir exploré l'implication ($\to$), notre voyage au cœur de la correspondance de Curry-Howard se poursuit. Aujourd'hui, nous allons découvrir comment Rocq traite la Conjonction (ET), la Disjonction (OU), et l'Absurdité (NON).

Nous adopterons une perspective profondément constructiviste (la sémantique Brouwer-Heyting-Kolmogorov, ou BHK), où chaque connecteur logique n'est pas une simple table de vérité, mais une véritable \textbf{structure de données} informatique construite par nos soins !

\begin{lstlisting}
Notation "( x , y )" := (conj x y).
Notation "'ι₁'" := or_introl (at level 0).
Notation "'ι₂'" := or_intror (at level 0).

Section Lesson2.

Variables A B C : Prop.
\end{lstlisting}

\section{La Conjonction (Le ET : $\land$)}

En logique de la réalisation (constructivisme), posséder une preuve formelle de $A \land B$, c'est posséder indépendamment une preuve de $A$ ET une preuve de $B$. Informatiquement, cela correspond exactement à la notion de \textbf{paire} (ou ``tuple'').

La règle d'introduction de la conjonction en déduction naturelle est la suivante :
$$\prftree[r]{${\land}I$}{\Gamma \vdash M : A}{\Gamma \vdash N : B}{\Gamma \vdash (M, N) : A \land B}$$

Pour construire/introduire une telle paire dans Rocq, nous utilisons la tactique \texttt{split}.

\begin{lstlisting}
Lemma ET_intro : A -> B -> A /\ B.
Proof.
  intros x y.
  split.
  - exact x.
  - exact y.
Qed.

Print ET_intro.
\end{lstlisting}

\begin{remarque}{}{}
  Observons le terme final. Le compilateur nous montre que le terme de preuve est bel et bien structuré comme une paire !
\end{remarque}

À l'inverse, si nous \emph{avons} une hypothèse de type $A \land B$, nous pouvons la ``détruire'' pour en extraire les deux composants sous-jacents, grâce à la commande \texttt{destruct}. Cela correspond aux \textbf{projections} $\pi_1$ et $\pi_2$ :
$$\prftree[r]{${\land}E_1$}{\Gamma \vdash M : A \land B}{\Gamma \vdash \pi_1 M : A} \qquad \prftree[r]{${\land}E_2$}{\Gamma \vdash M : A \land B}{\Gamma \vdash \pi_2 M : B}$$

\begin{lstlisting}
Lemma ET_symetrie : A /\ B -> B /\ A.
Proof.
  intros H_ab.
  destruct H_ab as [x y].
  split.
  - exact y.
  - exact x.
Qed.
\end{lstlisting}

\section{La Disjonction (Le OU : $\lor$)}

Voici le grand point de rupture du constructivisme ! Pour prouver $A \lor B$, vous ne pouvez pas joyeusement affirmer ``bon, l'un des deux doit être vrai'', sans savoir lequel. Vous DEVEZ exhiber consciemment une preuve ou bien de $A$, ou bien de $B$.

Informatiquement, c'est ce qu'on appelle un \textbf{type somme}. L'ordinateur conserve la trace indélébile du choix effectué via des fonctions d'injection : \texttt{left} correspond à $\iota_1$ et \texttt{right} correspond à $\iota_2$.

$$\prftree[r]{${\lor}I_1$}{\Gamma \vdash M : A}{\Gamma \vdash \iota_1 M : A \lor B} \qquad \prftree[r]{${\lor}I_2$}{\Gamma \vdash N : B}{\Gamma \vdash \iota_2 N : A \lor B}$$

\begin{lstlisting}
Lemma OU_intro_gauche : A -> A \/ B.
Proof.
  intros x.
  left.
  exact x.
Qed.

(** On constate que le terme de preuve final est l'application de l'injection 
    gauche [ι₁] à la variable [x] de type [A].
    On note aussi que le type [B] est passé en argument juste avant [x] : c'est
    parce que Rocq a besoin de stocker le type de la branche inutilisée... *)
Print OU_intro_gauche.
\end{lstlisting}

\begin{remarque}{Structures de Données vs Contrôle}{val}
  En informatique, on sépare généralement deux grands piliers :
  \begin{enumerate}
    \item \textbf{Structures de Données} : comment la machine représente l'information (ex: les paires pour le ET, ou les injections pour le OU).
    \item \textbf{Structures de Contrôle} : comment la machine prend des décisions et oriente le flux du calcul (ex: le \texttt{if... then... else...}).
  \end{enumerate}

  Posséder une disjonction $A \lor B$ en hypothèse forme une structure de donnée \textbf{fermée}, car on ne sait pas a priori quel choix a été fait (injection gauche ou droite). Pour l'exploiter, nous avons besoin d'une structure de contrôle pour ``ouvrir'' cette boîte. La généralisation du \texttt{if} s'appelle le \textbf{Filtrage par Motif} (\texttt{match ... with}). 
  
  C'est un aiguillage : ``SI la boîte contient une preuve de $A$, ALORS applique ce raisonnement ; SI elle contient une preuve de $B$, ALORS applique cet autre raisonnement''. C'est ce que génère la tactique \texttt{destruct}.
\end{remarque}

\begin{lstlisting}
Lemma OU_symetrie : A \/ B -> B \/ A.
Proof.
  intros H_ou.
  destruct H_ou as [x | y].
  - (* Cas 1 : La preuve d'origine etait une injection gauche, pour A *)
    right. exact x.
  - (* Cas 2 : La preuve d'origine etait une injection droite, pour B *)
    left. exact y.
Qed.
\end{lstlisting}

\section{L'Absurdité et la Négation}

Dans un monde constructiviste, la définition de la négation est d'une élégance presque poétique : $\neg A$ est exactement une \textbf{abréviation} pour $A \to \bot$ (où $\bot$ dénote l'absurdité, \texttt{False} dans Rocq). Prouver que $A$ est faux, c'est écrire une fonction qui, si elle recevait une preuve de $A$, réussirait à produire l'absurdité.

Et qu'est-ce que \texttt{False} ? C'est le type ultimement vide : il n'a \textbf{aucun habitant}.

\begin{lstlisting}
Lemma non_contradiction : ~(A /\ ~A).
Proof.
  intros H_impossible.
  destruct H_impossible as [a non_a].
  apply non_a.
  exact a.
Qed.
\end{lstlisting}

\subsection*{Le Principe d'Explosion (\emph{Ex Falso Quodlibet})}

Si vous obtenez une preuve de \texttt{False} dans votre contexte, c'est que vous avez une contradiction. Et d'une contradiction, vous pouvez déduire N'IMPORTE QUOI. En Rocq, c'est la tactique \texttt{exfalso}.

$$\prftree[r]{$\bot E$}{\Gamma \vdash M : \bot}{\Gamma \vdash \text{abort } M : A}$$

\begin{lstlisting}
Lemma explosion : False -> A.
Proof.
  intro f.
  exfalso.
  exact f.
Qed.
\end{lstlisting}

\section{La Crise du Constructivisme et le Tiers-Exclu}

L'idéalisme intuitionniste a un revers : le Tiers-Exclu $A \lor \neg A$ n'est pas prouvable de manière générique. Réussir à le prouver reviendrait à posséder l'Algorithme Universel capable de trancher la vérité de n'importe quelle proposition sans calcul. Ce rêve de David Hilbert a été ruiné par Gödel et Turing.

\begin{lstlisting}
Lemma tiers_exclu_fail : A \/ ~A.
Proof.
  left. Undo.
  right. intro H.
Abort.
\end{lstlisting}

Pour retrouver la logique classique, on doit imposer le Tiers-Exclu comme une vérité infuse via la commande \texttt{Axiom}.

\begin{lstlisting}
Axiom LEM : forall X : Prop, X \/ ~X.

Lemma double_negation : ~~A -> A.
Proof.
  intros non_non_a.
  destruct (LEM A) as [x | non_a].
  - exact x.
  - exfalso. apply non_non_a. exact non_a.
Qed.
\end{lstlisting}

\section{Exercices à faire}

\begin{exercice}{Lois de De Morgan (Sens Intuitionniste)}{}
  Prouver formellement $(A \lor B) \to \neg(\neg A \land \neg B)$. Ce sens est constructif, nul besoin du Tiers-Exclu !
\end{exercice}

\begin{exercice}{Lois de De Morgan (Sens Classique)}{}
  Prouver la réciproque $\neg(\neg A \land \neg B) \to (A \lor B)$. Attention : ici le constructivisme est paralysé, vous devrez utiliser l'axiome \texttt{LEM}.
\end{exercice}

\begin{lstlisting}
End Lesson2.
\end{lstlisting}

\end{document}
